%%
%% fall-lecture13.tex
%% 
%% Made by Alex Nelson <pqnelson@gmail.com>
%% Login   <alex@lisp>
%% 
%% Started on  2026-10-27T10:14:49-0700
%% Last update 2026-10-27T10:14:49-0700
%% 

\lecture[Hilbert spaces]{}

\begin{definition}
A \define{Pre-Hilbert Space} is a vector space $X$ over $\FF$ (here
$\FF$ is either the reals or
complex numbers) equipped with an inner product $(-,-)\colon X\times X\to\FF$,
which is a form satisfying
\begin{enumerate}
\item it is linear in the first slot: $(x_{1}+x_{2},y)=(x_{1},y)+(x_{2},y)$
and $(\lambda x,y)=\lambda(x,y)$
\item it is strictly positive: $(x,x)\geq0$ and $(x,x)=0\iff x=0$
\item conjugate symmetry: $(x,y)=\overline{(y,x)}$.
\end{enumerate}
\end{definition}

\begin{definition}
For any pre-Hilbert space $X$, we define the \define{Induced Norm}
$\|-\|\colon X\to\RR$ by $\|x\|:=\sqrt{(x,x)}\geq0$.

We will later prove it is, in fact, a norm. But we will stipulate this
is the case, and that we endow the pre-Hilbert space with the induced
metric and associated metric topology.
\end{definition}

\begin{fact}[Cauchy--Schwarz inequality]
For any pre-Hilbert space, we have
\begin{equation}
|(x,y)|\leq\|x\|\cdot\|y\|.
\end{equation}
The proof is tedious and not terribly enlightening.
\end{fact}

\begin{fact}[Parallelogram law]
For any pre-Hilbert space $X$ and for any $x,y\in X$, we have
\begin{equation}
\|x+y\|^{2}+\|x-y\|^{2}=2(\|x\|^{2}+\|y\|^{2}).
\end{equation}
This follows by rewriting everything with the inner product.
\end{fact}

\begin{fact}[Polarization identity]
\begin{subequations}
  \begin{enumerate}
\item  If $X$ is a real normed vector space, then
\begin{equation}
2(x,y) = \|x+y\|^{2}-\|x\|^{2}-\|y\|^{2}.
\end{equation}
\item If $X$ is a complex normed vector space, then
  \begin{equation}
4(x,y) = \|x+y\|^{2}-\|x-y\|^{2}+\I\|x+\I y\|^{2}-\I\|x-\I y\|^{2}.
  \end{equation}
\end{enumerate}
\end{subequations}
\end{fact}

\begin{theorem}
\begin{enumerate}
\item The induced norm on a pre-Hilbert space really is a norm.
\item Normed vector spaces are pre-Hilbert if and only if it satisfies
  the parallelogram identity.
\item The inner product is continuous: if $(x_{n})\to x$ and
  $(y_{n})\to y$, then $(x_{n},y_{n})\to (x,y)$.
\end{enumerate}
\end{theorem}

\begin{proof}
\begin{enumerate}
\item Follows from the Cauchy--Schwarz inequality.
\item Use polarization identity to obtain the candidate for the inner
  product and show it works fine.
\item By Cauchy--Schwarz inequality.\qedhere
\end{enumerate}
\end{proof}

\begin{corollary}
For any subset $A$ of a pre-Hilbert space $X$, we define
$A^{\perp}:=\{x\in X\mid\forall a\in A\ldotp (x,a)=0\}$.
Then $A^{\perp}$ is a closed subspace of $X$.
\end{corollary}

\begin{remark}
It is common to just use the term ``subspace'' for \emph{closed}
subspaces in the theory of Banach spaces and Hilbert spaces, because
closed subspaces are the interesting guys here.
\end{remark}

\begin{definition}
A \define{Hilbert Space} is a complete pre-Hilbert space. (Remember:
we endow pre-Hilbert spaces with the additional structure of a normed
vector space, using the induced norm, and thereby give it a metric and
the metric topology.)
\end{definition}

\begin{example}
The space $\ell^{2}(\NN)$ of square summable (real or complex)
sequences with norm
\begin{equation}
\|x_{n}\|^{2}=\sum_{n}|x_{n}|^{2}.
\end{equation}
We claim this norm satisfies the parallelogram identity, so it is a
pre-Hilbert space. But we know $\ell^{2}(\NN)$ is complete with
respect to the metric induced by the norm, so it is a Hilbert space.
\end{example}

\begin{non-example}
The space $\ell^{1}(\NN)$ of absolute summable (real or complex)
sequences with norm
\begin{equation}
\|x_{n}\|=\sum_{n}|x_{n}|.
\end{equation}
Although this is a Banach space, the parallelogram law is \emph{not}
satisfied. Therefore it is \emph{not} a Hilbert space.
\end{non-example}

\begin{lemma}[Riesz]
Let $H$ be a Hilbert space, let $A\subset H$ be a subset.
If $A$ is convex and closed, then there exists a unique $x\in A$ with
minimum norm.
\end{lemma}

\begin{proof}
\textsc{Step 1: Existence}. Define $d=\inf_{x\in A}\|x\|$. Let
$(x_{n})$ be a sequence in $A$ such that $\|x_{n}\|\to d$. Then
$(x_{n})$ is Cauchy. We see,
\begin{equation}\label{eq:math110a:lec13:eq1}
\|x_{n}-x_{m}\|^{2}=2(\|x_{n}\|^{2}+\|x_{m}\|^{2})-\|x_{n}+x_{m}\|^{2}
\end{equation}
by adding $\|x_{n}+x_{m}\|^{2}$ to both sides of the parallelogram
identity; and $2(\|x_{m}\|^{2}+\|x_{n}\|^{2})\to4d^{2}$ as $m,n\to\infty$;
and using convexity [$\forall a,b\in A\ldotp\frac{1}{2}(a+b)\in A$]
tells us
\begin{equation}
4\left\|\frac{x_{n}+x_{m}}{2}\right\|^{2}\geq4d^{2},
\end{equation}
since $x_{12}=(x_{1}+x_{2})/2\in A$ and by definition of $d$ we must
have $x_{12}$ have norm no smaller than $d$. Then we see the
right-hand side of Eq~\eqref{eq:math110a:lec13:eq1} becomes
\begin{equation}
4d^{2}-4\left\|\frac{x_{n}+x_{m}}{2}\right\|^{2}\leq0,
\end{equation}
whereas the left-hand side of Eq~\eqref{eq:math110a:lec13:eq1} is non-negative,
and therefore both must vanish. Hence $x_{n}$ must converge to a point
$x$ of norm $\|x\|=d$ (and by $A$ being closed, $x\in A$).

\textsc{Claim: Uniqueness}. If we have two candidates $\|x'\|=\|x''\|=d$,
then
\begin{subequations}
  \begin{align}
\|x'-x''\|^{2} &=
2(\|x'\|^{2}+\|x''\|^{2})-\|x'+x''\|^{2}\quad\mbox{by Eq~\eqref{eq:math110a:lec13:eq1}}\\
&=4d^{2}-\|x'+x''\|^{2},
  \end{align}
\end{subequations}
and by convexity $\|x'+x''\|^{2}\geq4d^{2}$, so this forces
\begin{equation}
\|x'-x''\|^{2}=0.
\end{equation}
Hence $x'=x''$.
\end{proof}

\begin{theorem}[Riesz projection]
Let $H$ be a Hilbert space, let $L\subset H$ be a closed subspace, let
$x\in H$ be a point. Then there exists a unique decomposition $x=x_{1}+y$
where $x_{1}\in L$ and $y\in L^{\perp}$ (i.e., $y\perp L$).
Equivalently, $H=L\oplus L^{\perp}$.
\end{theorem}

\begin{proof}
Take $x_{1}$ with minimal norm $\|x_{1}-x\|=d$ (apply Riesz Lemma to
the set $A=x+L$.) This defines $x_{1}$ in a unique way by the
Lemma. We just need to prove for any $z\in L$ we have $(x-x_{1},z)=0$.
We have, for any $t\in\CC$,
\begin{subequations}
  \begin{align}
d^{2} &\leq\|x-x_{1}+tz\|^{2}\\
&\leq (x-x_{1},x-x_{1}) + 2\Re(x-x_{1},tz)+(tz,tz).
  \end{align}
\end{subequations}
Then $2\Re(x-x_{1},tz)=0$ by positivity of $d^{2}$ (and the sum of a
linear and quadratic term must have the linear term vanish for them to
be positive). Hence the claim.
\end{proof}